MODELO MUNDIAL 2026 - COMPETENCIA DE MODELOS, CONTEXTO COMPETITIVO Y CLASIFICADOS

Este proyecto construye un sistema predictivo para estimar resultados de partidos del Mundial 2026.

El sistema entrena varios modelos de regresión, predice goles esperados para cada equipo, ajusta esas predicciones con contexto competitivo de fase de grupos y finalmente simula las tablas finales de grupo para determinar qué selecciones clasificarían a dieciseisavos de final.

============================================================
OBJETIVO DEL PROYECTO
============================================================

El objetivo principal es predecir:

- Marcadores probables de cada partido.
- Probabilidad de victoria local, empate y victoria visitante.
- Top 10 marcadores más probables.
- Modelo ganador entre varios algoritmos.
- Contexto competitivo de cada partido.
- Tablas finales predichas de cada grupo.
- Primeros y segundos de grupo.
- Mejores terceros clasificados.
- Posibles cruces de dieciseisavos de final.

============================================================
ESTRUCTURA ESPERADA
============================================================

modelo regresion/
│
├── data/
│   ├── results.csv
│   ├── results_backup_antes_mundial_2026.csv
│   ├── fifa_ranking.csv
│   ├── elo_ranking.csv
│   ├── fixtures_2026.csv
│   ├── fixtures_2026_group_stage.csv
│   ├── fair_play.csv
│   ├── data_modelo.csv
│   ├── fixtures_modelo.csv
│   └── fixtures_contexto_grupos.csv
│
├── outputs/
│   ├── competencia_modelos.csv
│   ├── mejor_modelo.txt
│   ├── modelo_home.pkl
│   ├── modelo_away.pkl
│   ├── predicciones_fixture.csv
│   ├── predicciones_todos_modelos.csv
│   ├── top10_marcadores.csv
│   ├── top10_todos_modelos.csv
│   ├── partidos_grupo_resultado_modelo.csv
│   ├── tabla_grupos_predicha.csv
│   ├── clasificados_grupos.csv
│   ├── mejores_terceros.csv
│   ├── clasificados_16avos.csv
│   ├── cruces_16avos_predichos.csv
│   ├── panel_tablas_grupos.png
│   ├── panel_clasificados_16avos.png
│   └── panel_cruces_16avos.png
│
├── 00_actualizar_results_mundial.py
├── 01_construir_data_modelo.py
├── 02_entrenar_predecir_modelos.py
├── 03_visualizar_partido_modelos.py
├── 04_mostrar_modelo_final.py
├── 05_crear_contexto_grupos.py
├── 06_predecir_clasificados.py
├── requirements.txt
└── README.txt

============================================================
REQUISITOS
============================================================

Instalar las librerías necesarias:

pip install -r requirements.txt

Contenido del archivo requirements.txt:

numpy
pandas
scipy
scikit-learn
matplotlib
joblib

============================================================
FLUJO DE EJECUCIÓN COMPLETO
============================================================

1) Entrar a la carpeta del proyecto:

cd "/home/jeanki/Escritorio/modelo regresion"

2) Actualizar resultados reales del Mundial:

python 00_actualizar_results_mundial.py

Este script actualiza:

data/results.csv

3) Construir la data del modelo:

python 01_construir_data_modelo.py

Genera:

data/data_modelo.csv
data/fixtures_modelo.csv

4) Crear contexto competitivo de fase de grupos:

python 05_crear_contexto_grupos.py

Genera:

data/fixtures_contexto_grupos.csv

5) Entrenar modelos y generar predicciones:

python 02_entrenar_predecir_modelos.py

Genera:

outputs/competencia_modelos.csv
outputs/mejor_modelo.txt
outputs/modelo_home.pkl
outputs/modelo_away.pkl
outputs/predicciones_fixture.csv
outputs/predicciones_todos_modelos.csv
outputs/top10_marcadores.csv
outputs/top10_todos_modelos.csv

6) Visualizar un partido específico:

python 03_visualizar_partido_modelos.py

Este script permite visualizar un partido por MATCH_ID.

Dentro del archivo se puede cambiar:

MATCH_ID = 60

El script muestra:

- partido seleccionado
- equipo local
- equipo visitante
- contexto competitivo
- comparación de predicciones por modelo
- top 10 marcadores por modelo
- gráfico comparativo de probabilidades
- gráfico pastel del modelo ganador
- tabla visual del contexto competitivo

7) Mostrar el modelo final:

python 04_mostrar_modelo_final.py

Este script muestra:

- modelo ganador
- coeficientes si el modelo es lineal
- importancia de variables si el modelo lo permite
- ecuación o interpretación del modelo final

8) Predecir clasificados a dieciseisavos:

python 06_predecir_clasificados.py

Este script construye la tabla final predicha de cada grupo y determina:

- primeros de grupo
- segundos de grupo
- mejores terceros
- clasificados a dieciseisavos
- cruces simulados de dieciseisavos

Genera:

outputs/partidos_grupo_resultado_modelo.csv
outputs/tabla_grupos_predicha.csv
outputs/clasificados_grupos.csv
outputs/mejores_terceros.csv
outputs/clasificados_16avos.csv
outputs/cruces_16avos_predichos.csv
outputs/panel_tablas_grupos.png
outputs/panel_clasificados_16avos.png
outputs/panel_cruces_16avos.png

============================================================
ORDEN CORRECTO DE EJECUCIÓN
============================================================

python 00_actualizar_results_mundial.py
python 01_construir_data_modelo.py
python 05_crear_contexto_grupos.py
python 02_entrenar_predecir_modelos.py
python 03_visualizar_partido_modelos.py
python 04_mostrar_modelo_final.py
python 06_predecir_clasificados.py

============================================================
DESCRIPCIÓN DEL MODELO
============================================================

El sistema entrena varios modelos para predecir goles esperados:

- Regresión Lineal
- Ridge
- Random Forest Regressor
- Gradient Boosting Regressor
- Poisson Regressor

Cada modelo se entrena dos veces:

- modelo_home: predice goles del equipo local
- modelo_away: predice goles del equipo visitante

Luego, los goles esperados se transforman en probabilidades usando una matriz de Poisson con ajuste Dixon-Coles.

============================================================
VARIABLES PRINCIPALES DEL MODELO BASE
============================================================

Las variables principales usadas para el entrenamiento son:

- home_gf12
- home_ga12
- home_pts12
- home_prev_matches
- away_gf12
- away_ga12
- away_pts12
- away_prev_matches
- diff_fifa
- diff_elo
- h2h
- neutral
- home_advantage

Estas variables representan forma reciente, goles, puntos, ranking FIFA, Elo, historial directo y localía.

============================================================
COMPETENCIA DE MODELOS
============================================================

El archivo 02_entrenar_predecir_modelos.py compara varios modelos usando métricas como:

- MAE
- RMSE
- R2
- Accuracy del resultado H/D/A

El mejor modelo se selecciona principalmente por el menor MAE promedio:

MAE promedio = (MAE local + MAE visitante) / 2

El modelo ganador se guarda en:

outputs/mejor_modelo.txt

============================================================
CONTEXTO COMPETITIVO
============================================================

El contexto competitivo se calcula en:

05_crear_contexto_grupos.py

Este contexto no se usa directamente en el entrenamiento porque la data histórica no contiene esas variables.

Se aplica como ajuste posterior sobre los goles esperados.

Primero, el modelo predice:

lambda_home_base
lambda_away_base

Luego el contexto ajusta:

lambda_home_final = lambda_home_base * factor_contexto_home
lambda_away_final = lambda_away_base * factor_contexto_away

Finalmente, con las lambdas ajustadas se calculan:

- probabilidad de victoria local
- probabilidad de empate
- probabilidad de victoria visitante
- top 10 marcadores más probables

============================================================
VARIABLES DE CONTEXTO COMPETITIVO
============================================================

El archivo fixtures_contexto_grupos.csv incluye variables como:

- puntos actuales del grupo
- posición actual
- mejor posición posible
- peor posición posible
- probabilidad de terminar en top 2
- probabilidad de terminar primero
- si el equipo necesita ganar
- si necesita resultado
- si pelea el primer lugar
- si puede caer a tercer lugar
- si necesita diferencia de goles
- si puede rotar
- nota interpretativa del contexto

Ejemplos de variables:

home_must_win
away_must_win
home_needs_result
away_needs_result
home_playing_for_first
away_playing_for_first
home_risk_drop_to_third
away_risk_drop_to_third
home_rotation_risk
away_rotation_risk

============================================================
INTERPRETACIÓN DEL FACTOR CONTEXTO
============================================================

Si un equipo necesita ganar, su factor contextual puede ser mayor que 1.

Ejemplo:

factor_contexto_home = 1.15

Esto significa que el contexto aumenta la intensidad ofensiva estimada del equipo local.

Si un equipo ya está clasificado y no pelea el primer lugar, puede tener:

factor_contexto_home = 0.90

Esto significa que el contexto reduce su intensidad ofensiva por posible rotación.

Si no hay ajuste contextual:

factor_contexto_home = 1.00

============================================================
CLASIFICACIÓN A DIECISEISAVOS
============================================================

El archivo 06_predecir_clasificados.py reconstruye la tabla final de cada grupo usando:

- resultados reales si ya existen en data/results.csv
- marcadores predichos si el partido aún no se jugó

Luego ordena cada grupo y clasifica:

- 12 primeros de grupo
- 12 segundos de grupo
- 8 mejores terceros

Total:

32 clasificados a dieciseisavos.

============================================================
CRITERIOS PARA ORDENAR GRUPOS Y MEJORES TERCEROS
============================================================

1. Puntos
2. Diferencia de goles
3. Goles a favor
4. Fair Play
5. Ranking FIFA

============================================================
FAIR PLAY
============================================================

El archivo opcional de Fair Play debe estar en:

data/fair_play.csv

Formato esperado:

team,yellow_cards,indirect_red_cards,direct_red_cards,team_conduct_score,source_note

Ejemplo:

Spain,2,0,0,-2,Team Conduct Score converted to yellow-card-equivalent counts for model tiebreaker
Uruguay,9,0,0,-9,Team Conduct Score converted to yellow-card-equivalent counts for model tiebreaker

El modelo calcula:

fair_play_score = yellow_cards*(-1) + indirect_red_cards*(-3) + direct_red_cards*(-4)

Si no existe data/fair_play.csv, el script asigna:

fair_play_score = 0

para todos los equipos.

============================================================
SALIDAS PRINCIPALES
============================================================

competencia_modelos.csv:
Contiene el rendimiento de cada modelo usando MAE, RMSE, R2 y accuracy del resultado.

mejor_modelo.txt:
Guarda el nombre del modelo ganador según MAE promedio.

predicciones_todos_modelos.csv:
Contiene predicciones de todos los modelos para todos los partidos.

predicciones_fixture.csv:
Contiene las predicciones del mejor modelo.

top10_todos_modelos.csv:
Contiene los 10 marcadores más probables por partido y por modelo.

top10_marcadores.csv:
Contiene los 10 marcadores más probables del mejor modelo.

partidos_grupo_resultado_modelo.csv:
Contiene los marcadores usados para construir las tablas. Puede incluir partidos reales y partidos predichos.

tabla_grupos_predicha.csv:
Contiene la tabla final predicha de cada grupo.

clasificados_grupos.csv:
Contiene primeros y segundos de cada grupo.

mejores_terceros.csv:
Contiene el ranking de terceros lugares y la columna advance.

clasificados_16avos.csv:
Contiene todos los clasificados a dieciseisavos.

cruces_16avos_predichos.csv:
Contiene una simulación de los cruces de dieciseisavos.

panel_tablas_grupos.png:
Imagen con las tablas finales predichas de todos los grupos.

panel_clasificados_16avos.png:
Imagen con primeros, segundos y mejores terceros.

panel_cruces_16avos.png:
Imagen con los cruces simulados de dieciseisavos.

============================================================
NOTA FINAL
============================================================

El modelo no garantiza el resultado real de los partidos.

Su objetivo es estimar probabilidades usando información histórica, forma reciente, ranking FIFA, Elo, historial, localía, contexto competitivo del grupo y reglas de clasificación.
